% Generated by support/export-artifacts.py; do not edit.
% Slug: ch61-provenance
% Source: chapters/ch61.tex:205
% TikZ source: chapters/ch61.tex:207
% Original label: \label{fig:ch61-provenance}
% Original caption (LaTeX source):
%   Retaining multiple origins supports explanation without defining a unique inverse edit.
% Caption: Retaining multiple origins supports explanation without defining a unique inverse edit.
% Setup: support/compiler-pseudocode.sty (TikZ libraries and shared styles).
\begin{tikzpicture}[node distance=15mm and 18mm]
\node[basic block] (a) {source: 2+3};
\node[basic block,right=of a] (b) {source: 1+4};
\node[basic block,below=of a] (c) {IR constant 5};
\draw[cfg edge] (a)--(c); \draw[cfg edge] (b)--(c);
\node[basic block,below=of c] (p) {origin set: both expressions};
\draw[cfg edge] (c)--(p);
\end{tikzpicture}

% Original figure LaTeX:
% \begin{figure}[htbp]
% \centering
% \begin{tikzpicture}[node distance=15mm and 18mm]
% \node[basic block] (a) {source: 2+3};
% \node[basic block,right=of a] (b) {source: 1+4};
% \node[basic block,below=of a] (c) {IR constant 5};
% \draw[cfg edge] (a)--(c); \draw[cfg edge] (b)--(c);
% \node[basic block,below=of c] (p) {origin set: both expressions};
% \draw[cfg edge] (c)--(p);
% \end{tikzpicture}
% \caption{Retaining multiple origins supports explanation without defining a unique inverse edit.}
% \label{fig:ch61-provenance}
% \end{figure}
